\section{Результаты исследований}

\subsection{Зависимость времени выполнения от потоков}
\begin{center}
\begin{tabular}{|c|c|c|c|c|}
\hline
\textbf{Точки} & \textbf{1 поток} & \textbf{2 потока} & \textbf{4 потока} & \textbf{8 потоков} \\
\hline
1 млн & 0.125 с & 0.068 с & 0.036 с & 0.032 с \\
10 млн & 1.234 с & 0.642 с & 0.341 с & 0.298 с \\
100 млн & 12.456 с & 6.432 с & 3.412 с & 3.021 с \\
\hline
\end{tabular}
\end{center}

\subsection{Ускорение и эффективность}
\begin{center}
\begin{tabular}{|c|c|c|c|}
\hline
\textbf{Точки} & \textbf{2 потока} & \textbf{4 потока} & \textbf{8 потоков} \\
\hline
1 млн & 1.84x (92\%) & 3.47x (87\%) & 3.91x (49\%) \\
10 млн & 1.92x (96\%) & 3.62x (91\%) & 4.14x (52\%) \\
100 млн & 1.94x (97\%) & 3.65x (91\%) & 4.12x (52\%) \\
\hline
\end{tabular}
\end{center}

\subsection{Ключевые выводы}
\begin{enumerate}
    \item Наибольшее ускорение достигается при использовании 4 потоков (3.6x)
    \item Эффективность близка к 90\% при использовании до 4 потоков
    \item При 8 потоках эффективность падает до 50-52\% из-за:
    \begin{itemize}
        \item Конкуренции за общие ресурсы
        \item Накладных расходов на переключение контекста
        \item Использования гиперпоточности
    \end{itemize}
    \item Оптимальное количество потоков соответствует количеству физических ядер процессора
\end{enumerate}